\documentclass[11pt,a4paper]{article}

% ---------------------------------------------------------
% PREAMBLE (stable, warning-free, Overleaf-safe)
% ---------------------------------------------------------
\usepackage[utf8]{inputenc}
\usepackage[T1]{fontenc}
\usepackage{lmodern}
\usepackage{amsmath, amssymb, amsfonts}
\usepackage{bm}
\usepackage{geometry}
\usepackage{graphicx}
\usepackage{hyperref}
\usepackage{cite}
\usepackage{authblk}
\usepackage{physics}
\usepackage{mathtools}

\geometry{margin=2.4cm}

\title{\textbf{Companion LII: The Lensing--ISW Cross-Correlation in the Relaxation-Driven Cyclic Cosmology}}
\author{Michael Lehmann \\ \texttt{mi.lehmann@gmx.de}}
\date{April 2026}

\begin{document}
\maketitle

\begin{abstract}
The cross-correlation between CMB lensing and the integrated Sachs--Wolfe (ISW) 
effect provides a sensitive probe of the temporal evolution of the gravitational 
potential. In the standard cosmological model, this correlation arises from the 
coupling between the time derivative of the potential and the integrated 
gravitational field along the line of sight. In the Relaxation-Driven Cyclic 
Cosmology (RDCC), the scale-dependent evolution of the gravitational potential 
modifies both the lensing kernel and the ISW source term, producing a distinctive 
signature in the lensing--ISW cross-correlation. This Companion develops a 
continuous description of the cross-correlation in RDCC, deriving how the 
relaxation scale influences the coupling between lensing and ISW, the coherence of 
large-scale modes, and the observational signatures in CMB--LSS analyses. The 
resulting framework provides a direct probe of the temporal structure of the RDCC 
gravitational field.
\end{abstract}
% ---------------------------------------------------------
\section{Introduction}

The cross-correlation between CMB lensing and the integrated Sachs--Wolfe (ISW) 
effect probes the coupling between the spatial and temporal structure of the 
gravitational potential. CMB lensing is sensitive to the integrated potential 
$\Phi + \Psi$, while the ISW effect is sensitive to the time derivative 
$\dot{\Phi} + \dot{\Psi}$. Their cross-correlation therefore encodes the 
relationship between the potential and its time evolution.

In the standard cosmological model, this correlation is determined by the onset of 
dark energy domination. In RDCC, the gravitational potential evolves differently. 
The relaxation mechanism suppresses the decay of small-scale modes and enhances the 
coherence of large-scale modes. This modifies both the lensing potential and the 
ISW source term in a scale-dependent manner, producing a distinctive signature in 
the lensing--ISW cross-correlation.
% ---------------------------------------------------------
\section{The RDCC Lensing Potential}

The CMB lensing potential is given by
\begin{equation}
    \phi(\hat{\mathbf{n}})
    = -2 \int_{0}^{\chi_{*}}
      d\chi\,
      \frac{\chi_{*} - \chi}{\chi_{*}\chi}\,
      \left[\Phi + \Psi\right](\chi,\hat{\mathbf{n}}).
\end{equation}

In RDCC, the metric potentials evolve as
\begin{align}
    \Phi_{\mathrm{RDCC}}(k,a)
    &= \Phi(k,a)
       \left[
         1 - \chi_{\Phi}
         \left(\frac{k}{k_{\mathrm{rel}}}\right)^{\alpha}
       \right], \\
    \Psi_{\mathrm{RDCC}}(k,a)
    &= \Psi(k,a)
       \left[
         1 - \chi_{\Psi}
         \left(\frac{k}{k_{\mathrm{rel}}}\right)^{\alpha}
       \right].
\end{align}

The RDCC lensing potential becomes
\begin{equation}
    \phi_{\mathrm{RDCC}}
    = \phi
      \left[
        1 - \chi_{\phi}
        \left(\frac{k}{k_{\mathrm{rel}}}\right)^{\alpha}
      \right].
\end{equation}
% ---------------------------------------------------------
\section{The RDCC ISW Source Term}

The ISW temperature shift is given by
\begin{equation}
    \frac{\Delta T_{\mathrm{ISW}}}{T_{\gamma}}
    = 2 \int \dot{\Phi}\, d\chi.
\end{equation}

In RDCC, the time derivative of the potential becomes
\begin{equation}
    \dot{\Phi}_{\mathrm{RDCC}}(k,a)
    = \dot{\Phi}(k,a)
      \left[
        1 - \chi_{\Phi}
        \left(\frac{k}{k_{\mathrm{rel}}}\right)^{\alpha}
      \right]
      - \Phi(k,a)\,
        \chi_{\Phi}\,\alpha\,
        \left(\frac{k}{k_{\mathrm{rel}}}\right)^{\alpha}
        \frac{\dot{k}_{\mathrm{rel}}}{k_{\mathrm{rel}}}.
\end{equation}

The second term is unique to RDCC and provides a direct probe of the relaxation 
mechanism.
% ---------------------------------------------------------
\section{The Lensing--ISW Cross-Correlation}

The cross-correlation is defined as
\begin{equation}
    C_{L}^{\phi T}
    = \langle \phi_{LM} T^{*}_{LM} \rangle.
\end{equation}

In RDCC, this becomes
\begin{equation}
    C_{L}^{\phi T,\mathrm{RDCC}}
    = C_{L}^{\phi T}
      \left[
        1 - \chi_{\phi T}
        \left(\frac{L}{L_{\mathrm{rel}}}\right)^{\alpha}
      \right]
      + \Delta C_{L}^{\phi T,\mathrm{rel}},
\end{equation}
where the additional term is
\begin{equation}
    \Delta C_{L}^{\phi T,\mathrm{rel}}
    \propto
    \int d\chi\,
    \Phi(k,a)\,
    \chi_{\Phi}\,\alpha\,
    \left(\frac{k}{k_{\mathrm{rel}}}\right)^{\alpha}
    \frac{\dot{k}_{\mathrm{rel}}}{k_{\mathrm{rel}}}.
\end{equation}

This term is the hallmark of RDCC.
% ---------------------------------------------------------
\section{Large-Scale Coherence and RDCC Signatures}

The relaxation mechanism enhances the coherence of large-scale modes. The 
lensing--ISW cross-correlation therefore exhibits:

\begin{itemize}
    \item enhanced power at large angular scales,
    \item suppressed power at small scales,
    \item a characteristic turnover near $L_{\mathrm{rel}}$,
    \item a non-standard phase relation between lensing and ISW.
\end{itemize}

These signatures provide a direct observational probe of the temporal structure of 
the RDCC gravitational field.
% ---------------------------------------------------------
\section{Conclusion}

The lensing--ISW cross-correlation in the Relaxation-Driven Cyclic Cosmology 
reflects the scale-dependent evolution of the gravitational potential. The 
relaxation mechanism modifies both the lensing potential and the ISW source term in 
a scale-dependent manner, producing distinctive signatures in the cross-correlation 
spectrum. These effects provide a direct observational window into the relaxation 
scale and the temporal structure of the RDCC gravitational field. The present 
Companion establishes the theoretical foundation for future studies of CMB--LSS 
coupling, potential evolution, and the interplay between matter and gravity in 
RDCC.

\bibliographystyle{apsrev4-2}
\nocite{*}
\bibliography{references}

\end{document}
